\documentclass{article}


% if you need to pass options to natbib, use, e.g.:
%     \PassOptionsToPackage{numbers, compress}{natbib}
% before loading neurips_2024


% ready for submission
% \usepackage[final]{neurips_2024}


% to compile a preprint version, e.g., for submission to arXiv, add add the
% [preprint] option:
%     \usepackage[preprint]{neurips_2024}


% to compile a camera-ready version, add the [final] option, e.g.:
%     \usepackage[final]{neurips_2024}


% to avoid loading the natbib package, add option nonatbib:
%    \usepackage[nonatbib]{neurips_2024}


\usepackage[utf8]{inputenc} % allow utf-8 input
\usepackage[T1]{fontenc}    % use 8-bit T1 fonts
\usepackage{hyperref}       % hyperlinks
\usepackage{url}            % simple URL typesetting
\usepackage{booktabs}       % professional-quality tables
\usepackage{amsfonts}       % blackboard math symbols
\usepackage{nicefrac}       % compact symbols for 1/2, etc.
\usepackage{microtype}      % microtypography
\usepackage{xcolor}         % colors
\usepackage{graphicx}       % figures
\graphicspath{{./}{preliminary_report/}}


\title{Preliminary Report For Project 7\\Sentiment Analysis}


% The \author macro works with any number of authors. There are two commands
% used to separate the names and addresses of multiple authors: \And and \AND.
%
% Using \And between authors leaves it to LaTeX to determine where to break the
% lines. Using \AND forces a line break at that point. So, if LaTeX puts 3 of 4
% authors names on the first line, and the last on the second line, try using
% \AND instead of \And before the third author name.


\author{
  Group 19 \\
  Nolan Capomaggio and Mathis Morra-Fischer \\
%   \texttt{\url{https://github.com/orion6002/EARIN\_lab\_group\_19}}
}


\begin{document}


\maketitle

\begin{abstract}
  Sentiment analysis is a key task in Natural Language Processing (NLP), 
  aiming to automatically determine the emotional polarity of a given text. 
  In this work, the goal is to address binary sentiment classification of Amazon 
  product reviews. We distinguish positive (4–5 stars) from negative (1–2 stars) 
  review. We use the Amazon Reviews dataset (Kaggle, Adam Bittlingmayer), 
  which includes approximately 4 million labeled reviews. We propose to compare 
  three approaches of increasing complexity: Logistic Regression with TF-IDF 
  features, a LSTM neural network, and a fine-tuned RoBERTa 
  model. Models are evaluated using accuracy and macro F1-score.
\end{abstract}


\section{Introduction}
For large companies as Amazon, being able to classify good and bad reviews is 
mandatory in order to be able to satisfy the custommer. Knowing this, our 
goal is to developp a thrustworthy solution to tackle this issue. Our solution 
is known as sentiment analysis.

In this project, we frame sentiment analysis as a \textbf{binary text classification}
problem.
Given the text of an Amazon product review, the goal is to predict a label
$y \in \{0, 1\}$, where $0$ denotes a negative review (1–2 stars) and $1$ denotes
a positive review (4–5 stars).
Reviews rated 3 stars (neutral) are excluded from the dataset by construction,
so no neutral class is considered.

We use the Amazon Reviews dataset published by Bittlingmayer on Kaggle,
which provides 3.6 million training examples and 400,000 test examples in fastText
format.
We propose to compare three methods of increasing complexity: a Logistic Regression
classifier with TF-IDF features, a Long Short-Term Memory (LSTM) recurrent neural
network, and a fine-tuned RoBERTa transformer model.


\section{Dataset}
We use the Amazon Reviews dataset from Kaggle (\url{https://www.kaggle.com/datasets
/bittlingmayer/amazonreviews}). 
Each line is formatted as a fastText label followed by the review text:
\texttt{\_\_label\_\_1} corresponds to a negative review (1--2 stars) and
\texttt{\_\_label\_\_2} corresponds to a positive review (4--5 stars).
Three-star reviews are excluded by construction, so the task is strictly binary.
It contains approximately 4 million labeled reviews.

\subsection{Data Balance \& Splits}
The two classes are distributed equally: exactly 50\% negative and 50\% positive
in both the training and test splits.
The dataset is therefore \textbf{perfectly balanced}, and no oversampling or
undersampling is required.\\
The dataset contains approximately 4 millions entries. It is already splitted into 
a training part and a test part with a ratio of 90/10 (which is enough because the 
testing part is already consituated of 400 000 entries). This is considered a 
standard for this kind of dataset.

\subsection{Preprocessing}
For all methods implemented we have to do a based preprocess. The \texttt{\_\_label\_\_X}
prefix is stripped.
The text itself requires cleaning because it is written by humans which are not reliable 
and use special characters. We will remove HTML tags, special characters and unnecessary 
punctuation. We will also convert text to lowercase to prevent taking uppercased words 
as different words.\\
For each models that we are going to use, we have to do a different preprocessing in 
order to make them more efficient.\\

- For Logistic Regression we will use TF-IDF (Term Frequency-Inverse Document Frequency). 
This turns text into a numerical matrix. It highlights words that are important to s
pecific reviews.\\
- For LSTM, the text will be converted into sequences of integers. We will use padding 
to ensure all inputs have the same length. We will also use pre-trained word 
embeddings to represent semantic meaning.\\
- The RoBERTa is the one with the most advanced preprocessing. We have to use a 
byte-level Byte-Pair Encoding tokenizer. This handles rare words by breaking them into 
smaller sub-units. Then We will add the <s> (start of string) and </s> (separator) 
tokens. The model uses the representation of the first token for classification. And 
finaly we have to create masks to tell the model which parts of the input are real text 
and which are only padding.


\subsection{Evaluation Metrics}
We want to collect the following metrics during tests:
\begin{itemize}
  \item \textbf{Accuracy}: the proportion of correctly classified examples.
    This is a valid primary metric here because the dataset is perfectly balanced.
  \item \textbf{Macro F1-score}: the unweighted mean of per-class F1 scores.
    It provides a more robust basis for comparing methods, particularly as a
    cross-benchmark reference.
  \item \textbf{Precision and Recall} per class: to identify asymmetric error
    patterns (e.g., a model biased toward predicting positive).
  \item \textbf{Confusion matrix}: for qualitative analysis of error types.
\end{itemize} We will compare our results to the 
top scores on the website. Based on the Kaggle leaderboard and existing literature for this 
dataset, the state-of-the-art accuracy is approximately 92–94\%. Our goal is to approach 
this upper bound using the RoBERTa model.

\section{Solution Description}

\subsection{Overview}

For our project we will be using three different methods:
\textit{Logistic Regression} represent documents as sparse
bag-of-words or TF-IDF vectors and train linear or kernel-based classifiers
\cite{pang2002thumbs}.
\textit{LSTM} model is a sequential structure of text and
capture order-dependent patterns; Long Short-Term Memory (LSTM) networks are the 
most used in this type of processing.
\textit{RoBERTa} such as BERT and its successors achieve state-of-the-art results 
combining large-scale unsupervised pre-training with task-specific fine-tuning.

The choice of method focuses on quantifing the performance gain at each level of
complexity and to evaluate whether the additional computational cost of
more advanced models is justified on this particular dataset.

\subsection{Logistic Regression with TF-IDF}

Logistic Regression trained on TF-IDF features serves as our baseline.
The method is fast, fully interpretable, and has been shown to yield competitive
accuracy on large-scale review datasets.

\paragraph{Feature extraction}
TF-IDF (Term Frequency – Inverse Document Frequency) assigns a weight to each
word in a document that reflects how informative that word is relative to the
entire corpus.
We can construct a TF-IDF matrix using unigrams and bigrams, retaining the top
100,000 features ranked by term frequency.
We fit a Logistic Regression classifier with $\ell_2$ regularization using the
\texttt{scikit-learn} implementation.

\paragraph{Motivation}
This method acts as a strong reference point: any more complex model must
surpass it to justify its additional cost in training time and computational
resources, it is a fast method (not time consuming) but not so efficient.

\subsection{Long Short-Term Memory Network}

LSTM networks are a class of recurrent neural
networks designed to capture long-range dependencies in sequential data.
Each LSTM cell maintains a memory state that is updated at every time step
through learned gating mechanisms (input, forget, and output gates), which
mitigate the vanishing gradient problem of vanilla RNNs.

\paragraph{Training.}
The model is trained with the Adam optimizer and
binary cross-entropy loss for up to 5 epochs with early stopping based on
validation F1-score for allowing training in a resonable time.

\paragraph{Motivation.}
Compared to the TF-IDF baseline, LSTM captures word order and local context.
It represents an intermediate complexity level and does not require a GPU
for inference, making it a practical middle choice.

\subsection{RoBERTa}

RoBERTa (Robustly Optimized BERT Pretraining Approach) 
is a transformer encoder pre-trained on a large
English text corpus.
It improves over the original BERT model by removing
the Next Sentence Prediction objective, and training with larger batches and 
more data. These modifications consistently improve downstream classification 
performance.

\paragraph{Architecture.}
We fine-tune the \texttt{roberta-base} checkpoint 
from the Hugging Face Transformers library.
A linear classification head is added on top of the representation of the
\texttt{<s>} (start-of-sequence) token, which aggregates contextual information
from the entire input.

\paragraph{Fine-tuning.}
Fine-tuning is performed for up to 3 epochs with a learning rate of
$2 \times 10^{-5}$, linear warmup over the first 10\% of steps, and weight
decay of 0.01.
The batch size is set to 32.

\paragraph{Motivation.}
RoBERTa represents the state-of-the-art family of methods for text
classification.
Its inclusion allows us to measure the benefit of large-scale pre-trained
contextual representations over task-specific architectures trained from
scratch.

\section{Midterm Solution}

\subsection{Implemented Progress}
At this stage of the project, we have implemented and evaluated one complete
method: \textbf{Logistic Regression with TF-IDF features}. This model is used as
the baseline because it is fast to train, easy to reproduce, and provides a
strong reference before implementing more expensive neural models.

The current implementation includes the following parts:
\begin{itemize}
  \item a loader for the fastText dataset format, converting
  \texttt{\_\_label\_\_1} to class 0 and \texttt{\_\_label\_\_2} to class 1;
  \item a stratified train/validation split that keeps the 50/50 class balance;
  \item basic preprocessing, including HTML tag removal, lowercasing, and
  whitespace normalization;
  \item a \texttt{scikit-learn} pipeline with TF-IDF vectorization followed by
  Logistic Regression;
  \item evaluation code for accuracy, macro F1-score, precision, recall,
  classification reports, and confusion matrices;
  \item saving the trained model as a \texttt{.pkl} file for later reuse.
\end{itemize}

The Logistic Regression model uses unigram and bigram TF-IDF features, a maximum
vocabulary size of 100,000 features, sublinear term-frequency scaling, and
$\ell_2$ regularization with $C=1.0$.

\subsection{Reproducibility}
The experiment is reproducible because the data split uses a fixed random seed.
The default seed is \textbf{42}, and it can be changed through the command line
argument \texttt{--seed}. The full baseline experiment was run with:

\begin{verbatim}
python3 sentiment/train.py \
  --model logistic \
  --train sentiment/data/train.ft.txt \
  --test sentiment/data/test.ft.txt \
  --output_dir sentiment/out/logistic_full
\end{verbatim}

The full Kaggle training file contains 3,600,000 reviews. We split it into
2,880,000 training examples and 720,000 validation examples. The official test
file contains 400,000 reviews.

\subsection{Intermediate Results}
Table~\ref{tab:midterm-logistic-results} presents the current results on the
full dataset. The validation and test results are very close, which suggests
that the model generalizes consistently and that the split is representative.

\begin{table}[h]
  \centering
  \caption{Midterm results for Logistic Regression with TF-IDF.}
  \label{tab:midterm-logistic-results}
  \begin{tabular}{lrrrrr}
    \toprule
    Split & Examples & Accuracy & Macro F1 & Precision & Recall \\
    \midrule
    Validation & 720,000 & 0.9413 & 0.9413 & 0.9413 & 0.9413 \\
    Test & 400,000 & 0.9407 & 0.9407 & 0.9407 & 0.9407 \\
    \bottomrule
  \end{tabular}
\end{table}

\begin{figure}[h]
  \centering
  \includegraphics[width=0.55\linewidth]{confusion_logistic.png}
  \caption{Confusion matrix for the Logistic Regression baseline on the full
  test set.}
  \label{fig:midterm-logistic-confusion}
\end{figure}

The model obtains 94.07\% accuracy and 94.07\% macro F1-score on the full test
set. Precision and recall are also balanced between the negative and positive
classes, both being approximately 0.94. This shows that the model does not have
a strong bias toward one class, which is expected because the dataset is
perfectly balanced.

\subsection{Challenges, Findings, and Corrections}
The first challenge was practical dataset handling. The Kaggle files are large
and must be extracted before training.

The main finding so far is that the baseline is stronger than expected. In the
initial plan, Logistic Regression was presented mainly as a fast but less
powerful baseline. However, it already reaches around 94\% test accuracy on the
full dataset. This means that LSTM and RoBERTa must be evaluated not only by
accuracy, but also by training time, computational cost, and practical benefit
over this strong baseline.

\subsection{Plan for the Final Report}
For the final report, we plan to complete the comparison between the three
selected methods: Logistic Regression, LSTM, and RoBERTa. The final experiments
should include:
\begin{itemize}
  \item accuracy, macro F1-score, precision, recall, and confusion matrices for
  all implemented models;
  \item comparison of training time and hardware requirements;
  \item analysis of whether the LSTM and RoBERTa models improve enough over
  Logistic Regression to justify their additional complexity;
  \item qualitative error analysis on selected misclassified reviews;
  \item final discussion of the best trade-off between performance,
  interpretability, and computational cost.
\end{itemize}


\section{Conclusion}
We have described the dataset, evaluation protocol, and three methods for
binary sentiment classification of Amazon product reviews.
The Amazon Reviews dataset is large, perfectly balanced, and requires only
minimal preprocessing; no class rebalancing technique is needed.
Our three methods – Logistic Regression with TF-IDF, LSTM, and RoBERTa – 
span the full range from a lightweight and interpretable baseline to a
state-of-the-art transformer model, enabling a controlled comparison of
performance versus computational cost.
In the next stage of the project, we will implement all three pipelines,
compare their accuracy and macro F1-score on the validation set, and
investigate the impact of key design choices such as vocabulary size,
sequence length, and fine-tuning strategy.

\section*{References}

{\small
[1] NeurIPS\ (2024) Call for Papers\\ \url{https://neurips.cc/Conferences/2024/CallForPapers}

[2] Adam Bittlingmayer\ (2020) Amazon Reviews for Sentiment Analysis\\ \url{https://www.kaggle.com/datasets/bittlingmayer/amazonreviews}
}

\end{document}
